\chapter{¿Cuánta metateoría hace falta?}
\label{cap:objeto-y-meta}

El capítulo~\ref{cap:kernel} dejó los dos cálculos y el puente entre ellos; el
capítulo~\ref{cap:qpp} dirá qué se afirma con ellos. Entre los dos queda una pregunta que casi
ningún libro de incompletitud hace explícita, y que aquí no se puede eludir porque la metateoría
es una máquina concreta: \emph{¿cuánta fuerza lógica estamos gastando para hablar de la teoría
objeto, y cuánta hace falta de verdad?}

La pregunta no es ociosa. Gödel escribió su metateoría en prosa, pero deliberadamente
\textbf{finitaria}, para que Hilbert la aceptara: no valía razonar sobre las pruebas con más
recursos que los que se estaban poniendo en duda. Aquí la metateoría es Lean~4, que es
enormemente más fuerte que eso. Si nadie mide cuánto se usa, el libro estaría demostrando la
incompletitud de una aritmética débil desde una torre de la que no sabe nada. Este capítulo mide.

\begin{leccion}[La asimetría que hay que decir en voz alta]
La metateoría \textbf{no tiene que ser más fuerte} que la teoría de la que habla. Sólo tiene que
poder hablar de \emph{sintaxis finita}: fórmulas, listas de fórmulas, y la comprobación mecánica de
que una lista es una derivación. Es una exigencia de \emph{expresividad}, no de fuerza. Esto es lo
que hace posible el proyecto entero — y lo que hace que la pregunta «¿cuánta metateoría?» tenga
una respuesta pequeña.
\end{leccion}

\section{Lean no es la aritmética de Peano, y por tres razones}

Conviene descartar primero la respuesta cómoda. Sería elegante que la metateoría de este proyecto
fuera exactamente PA —la aritmética de primer orden con inducción—: el libro se cerraría sobre sí
mismo. No lo es, y por bastante margen.

\textbf{Primero, la inducción de Lean no es un esquema.} La de PA sí: un axioma \emph{por cada}
fórmula del lenguaje aritmético, infinitas numerables, todas cuantificando sólo sobre números. Eso
es un \defterm{esquema de axiomas}\nivelterm{objeto}. El principio de inducción que Lean genera
para un tipo inductivo es \textbf{un solo enunciado} que cuantifica sobre un motivo
\ident{C : Formula → Sort u} — es decir, sobre \emph{todos} los predicados de fórmulas. Eso es
inducción de segundo orden, y no hay manera de escribirla como esquema numerable.

Y no es una diferencia de manual: el proyecto la usa así de verdad. La demostración de que el
código de una fórmula es igual a su numeral va por recursión meta sobre \ident{Formula} con el
motivo $\varphi \mapsto \Prf(\text{\ident{formCode}}\ \varphi \doteq \text{\ident{numeral}}
(\text{\ident{codeNat}}\ \varphi))$, y $\Prf$ es a su vez un \ident{Prop} inductivo — un motivo que
en PA no se puede ni enunciar sin haber aritmetizado antes $\Prf$.

\leanfrag{prf_formCode_numeral}

\textbf{Segundo, los universos.} Lean tiene una jerarquía infinita de universos con \ident{Prop}
impredicativo. Su teoría de tipos se interpreta en \textsc{zfc} más una cantidad numerable de
cardinales inaccesibles, y dentro de Lean puede construirse un modelo de \textsc{zfc}; luego Lean
demuestra $\mathrm{Con}(\textsc{zfc})$, y con ella $\mathrm{Con}(\mathrm{PA})$, que PA no
demuestra.\footnote{M.~Carneiro, \emph{The Type Theory of Lean}, tesis de máster, Carnegie Mellon
University, 2019. Es la referencia estándar para el modelo; el enunciado exacto conviene
consultarlo allí.}

\textbf{Tercero, los axiomas.} \ident{propext}, \ident{Classical.choice} y \ident{Quot.sound} están
en el \defterm{footprint}\nivelterm{proyecto} de prácticamente todo lo que este libro cita: son la
base sancionada del proyecto. Ninguno de los tres tiene contrapartida aritmética.

\section{Cuánto de esa fuerza se usa: dos niveles, y ninguno más}

Que Lean sea fuerte no dice nada sobre cuánto se gasta. Eso se mide, y el resultado es
sorprendentemente pequeño.

\begin{center}\small
\begin{tabular}{lll}
\toprule
qué & tipo & nivel \\
\midrule
\ident{Term}, \ident{Formula}, \ident{Pos} & \ident{Type} & \ident{Sort 1} \\
\ident{Derives}, \ident{Prf}, \ident{PrfH} & \ident{Prop} & \ident{Sort 0} \\
\bottomrule
\end{tabular}
\end{center}

En todo el árbol de \modulo{ROBINSON_PlusPlus/} y en los módulos de \modulo{FOL} que la cadena de
\ident{import} alcanza no hay \textbf{ni un} \ident{Sort u}, \textbf{ni un} \ident{Type 1},
\textbf{ni una} declaración \ident{universe}. El único \ident{Type u} del árbol vive en
\modulo{FOL/Completeness.lean}, que —como dijo el capítulo~\ref{cap:kernel}— ningún módulo importa.

Aunque el recursor de \ident{Formula} sea polimórfico en universos, el proyecto lo instancia en
exactamente dos sitios: motivo en \ident{Prop} cuando \emph{demuestra} por inducción, y motivo en
\ident{Type} cuando \emph{define} por recursión (\ident{codeNat}, \ident{formCode},
\ident{liftTerm}, \ident{substFormula}). Ninguna inducción concreta necesita el polimorfismo: cada
una vive en un nivel fijo.

\begin{observacion}
Y algo más fuerte: \textbf{no se cuantifica sobre \ident{Prop} en ninguna parte} —ni
\ident{(P : Prop)} ni \ident{{P : Prop}}—, luego la impredicatividad de \ident{Prop}, que es de
donde sale buena parte de la fuerza de Lean, \textbf{no se usa}. Dato relacionado: hay un solo
\ident{termination_by} en todo el árbol (\modulo{ROBINSON_PlusPlus/Meta/CodeDecode.lean}); todo lo
demás es recursión estructural.
\end{observacion}

\begin{muro}[Qué mide exactamente esto]
Todas las cifras de esta sección están medidas sobre el \textbf{texto} del árbol, con búsquedas
literales. Un \printaxioms\ sobre el entorno compilado puede destapar usos implícitos que una
búsqueda no ve —una táctica que instancia algo en un universo alto, por ejemplo—. El resultado hay
que leerlo como «no hay ningún uso \emph{escrito}», que es menos de lo que uno querría y bastante
más de lo que suele decirse.
\end{muro}

\section{La traducción a una metateoría aritmética}

Dos niveles, y cada uno tiene contrapartida conocida en aritmética:

\begin{itemize}
\item motivo en \ident{Type} —definir una función recursiva sobre sintaxis— se traduce por
  \defterm{recursión de curso de valores}\nivelterm{meta}: la función $\beta$ de Gödel, que permite
  codificar una sucesión finita en un número y así definir por recursión dentro de la aritmética.
  En este proyecto la capa \ident{lenc}/\ident{nthc} ya se construye \emph{para otra cosa}, y es
  exactamente ese instrumental;
\item motivo en \ident{Prop} —probar por inducción— se traduce por el esquema de inducción de la
  metateoría.
\end{itemize}

La fuerza necesaria es del orden de $\mathrm{I}\Sigma_1$: inducción restringida a fórmulas
$\Sigma_1$. Y aquí se cierra un círculo que merece decirse: \doc{AXIOMS.md}~§1.1 argumenta que la
\emph{teoría objeto} necesita inducción de ese mismo perfil para D2 y D3. Teoría y metateoría
piden lo mismo. No es casualidad: las dos tienen que hacer el mismo trabajo, hablar de objetos
sintácticos finitos.

\section{\texorpdfstring{\texttt{Classical}}{Classical}: dos usos, y no son de la misma clase}

En todo \modulo{ROBINSON_PlusPlus/} hay \textbf{dos} usos explícitos de lógica clásica. Dos. Y la
diferencia entre ellos es el resultado más interesante de este capítulo.

\textbf{(a) Eliminable.} En \modulo{ROBINSON_PlusPlus/Full/PrimeFactor.lean} se saca
$\exists a,\ a \mid n \wedge a \neq 1 \wedge a \neq n$ por reducción al absurdo desde
«$n$ no es primo».

\leanfrag{exists_prime_factor}

Pero todos los divisores están acotados por $n$, y la divisibilidad, la igualdad y «ser primo» son
decidibles sobre \ident{Nat}. Es una \textbf{búsqueda acotada}: constructiva de manual.
\ident{Decidable.byContradiction} la sustituye sin perder nada.

\textbf{(b) No eliminable — y está identificado.} El otro está en la reducción que descarga la
hipótesis de reflexión:

\leanfrag{Reflects}
\leanfrag{reflects_of_omega}

Tras el \ident{intro}, el objetivo es $\Prf\varphi$ con la hipótesis $\neg\,\Prf\varphi$: o sea,
$\neg\neg\,\Prf\varphi \to \Prf\varphi$. Y $\Prf\varphi$ es $\Sigma_1$ —existe una derivación
finita, y comprobar que una lista dada es una derivación es decidible: ése es justo el trabajo del
verificador del capítulo~\ref{cap:verificador}—.

\begin{leccion}[Eso tiene nombre, y es el más débil que sirve]
Eliminar la doble negación sobre un enunciado $\Sigma_1$ de matriz decidible es el
\defterm{principio de Markov}\nivelterm{meta}. No es «lógica clásica» a secas: es el principio
\textbf{más débil} que separa la matemática constructiva intuicionista de la recursiva rusa. Poder
decir «un solo uso esencial, y es Markov» no es una nota de limpieza — es un resultado sobre el
proyecto.
\end{leccion}

El \ident{Classical.choice} que aparece en los footprints es otra cosa: entra casi todo por
instancias \ident{Decidable} y por tácticas. En
\modulo{ROBINSON_PlusPlus/Meta/ChainDecode.lean} se construyen a mano \ident{DecidableEq} para
\ident{Term} y \ident{Formula} —el \ident{deriving} automático no cubre el anidamiento
\ident{func : String → List Term}—, así que hay decidibilidad real disponible y ese uso es en
principio prescindible.

\begin{muro}[Lo que un \printaxioms\ no distingue]
Los footprints son la moneda del proyecto, pero \ident{Classical.choice} en un \printaxioms\ sale
igual venga del Markov esencial o de una táctica perezosa. Y esa distinción es justo la que hace
falta. La sección siguiente cuenta que la herramienta que la hace ya existe.
\end{muro}

\section{El experimento ya está hecho, y en el mismo ecosistema}

Sustituir la maquinaria meta de Lean por una metateoría aritmética es un objetivo declarado del
autor: usar su proyecto \modulo{Peano} —con alguna extensión— en lugar de Lean. La pregunta
natural es si eso es una aspiración o un plan. Resulta que es un plan, porque la parte difícil ya
se hizo allí.

\modulo{Peano} lleva desde 2026 una \textbf{guarda de compilación}: un comando que, en cada
\ident{lake} \ident{build}, comprueba que los teoremas vigilados no dependen de
\ident{Classical.choice}. Si alguien introduce lógica clásica por accidente, la compilación falla.

\leanfrag{peano_assert_constructive}

Hoy hay \textbf{1425} invocaciones activas de ese comando. Y —el detalle que lo hace confiable— el
fichero conserva, comentada, una invocación sobre un teorema que \emph{sí} usa elección, para poder
verificar que la guarda detecta lo que dice detectar. Es un control positivo, exactamente la
disciplina que este libro se impone a sí mismo en \doc{PLAN-LIBRO.md}~§2.

\begin{leccion}[Lo que eso demuestra para el plan de sustitución]
La técnica es \textbf{portable tal cual}: es un fichero, no una metodología por inventar.
Aplicarla aquí daría la separación que la sección anterior pide —Markov esencial frente a
\ident{Decidable} incidental— como \emph{control de compilación} y no como auditoría manual. Es la
doctrina de este libro aplicada a sí mismo: convertir un principio en una máquina.
\end{leccion}

Y hay un dato que vale más que el método. La función $\beta$ de Gödel —la pieza que la sección
anterior identifica como necesaria para que una metateoría aritmética hospede la recursión de curso
de valores— figuraba en la lista inicial de deuda clásica de \modulo{Peano}. Ya no: sus doce
declaraciones principales están hoy bajo la guarda, y la reparación consistió en reemplazar la
elección por \textbf{búsqueda acotada} — la misma técnica que arriba se propone para el uso
eliminable de este proyecto. La única excepción que \modulo{Peano} mantiene, documentada y
aceptada, es un módulo cuyo cometido es cuantificar sobre sistemas de Peano \emph{abstractos}
cualesquiera: metateoría genuinamente no constructiva, sin nada que enumerar ni acotar.

\begin{observacion}
Nótese qué acaba de pasar aquí, porque es un caso de \doc{DOCSTRINGS-NO-FIABLES.md} en vivo. La
lista de deuda clásica que la documentación de \modulo{Peano} exhibe sigue nombrando el módulo de
la función $\beta$; el árbol dice que se resolvió. La documentación es el inventario de partida, no
el estado. Lo que decide es la medición.
\end{observacion}

\section{Qué transferiría y qué no}

Suponiendo hecha la sustitución, conviene decir por adelantado cuál sería su alcance, porque no es
total y la mitad que falta no es un defecto.

\begin{center}\small
\begin{tabularx}{\linewidth}{@{}>{\raggedright\arraybackslash}p{0.22\linewidth}
                              >{\raggedright\arraybackslash}X@{}}
\toprule
qué & transfiere a una metateoría de fuerza $\mathrm{I}\Sigma_1$ \\
\midrule
todo $\Prf$ & \textbf{sí, entero}: es finitario, r.e.\ y por recursión estructural sobre objetos
finitos. D1, D2, el punto fijo y Gödel~II sobre Q\texttt{++} caben ahí \\
$\derives$ con \ident{gen} & \textbf{no}: la \defnot{$\omega$-regla} tiene premisas infinitas, y
ninguna aritmética de primer orden puede hospedarla \\
\bottomrule
\end{tabularx}
\end{center}

O sea: una metateoría-Peano reproduciría la mitad $\Prf$ del proyecto y no la mitad $\derives$. No
es una limitación nueva — es exactamente la misma razón por la que el capítulo~\ref{cap:kernel}
tenía que presentar \emph{dos} cálculos y no uno. Pero la sustitución habría que enunciarla como
parcial, y decir cuál mitad, que es lo que casi nunca se dice.
